\chapter{الآثار الاجتماعية والاقتصادية (الأجراء، المقاولة، سوق الشغل)}

%% ─────────────────────────────────────────────────────────
%% مقدمة الفصل
%% ─────────────────────────────────────────────────────────
\section*{مقدمة الفصل}
\addcontentsline{toc}{section}{مقدمة الفصل}

شهدت المنظومة الإنتاجية وسوق الشغل بالمغرب تحولات جذرية فرضتها إكراهات العولمة الاقتصادية واحتدام التنافسية الدولية، مما دفع الفاعلين الاقتصاديين إلى البحث الحثيث عن آليات تدبيرية مرنة تُمكّن من التكيف اللحظي مع تقلبات الأسواق. وقد شكل الانتقال من نموذج علاقة الشغل الثنائية الكلاسيكية القائمة على الاستقرار والأمان المهني إلى نموذج العلاقات الثلاثية الأطراف (التي تجمع بين المقاولة المستعملة، ومقاولة التشغيل المؤقت أو مقاول المناولة، والأجير) المنعطف الأبرز في السياسة التشغيلية المعاصرة.

وإذا كان المشرع المغربي قد سعى من خلال مدونة الشغل (القانون رقم 65.99) إلى إيجاد نقطة توازن دقيقة بين متطلبات "المرونة الاقتصادية" للمقاولات من جهة، ومقتضيات "الأمان الاجتماعي" للأجراء من جهة ثانية فيما يُعرف اصطلاحاً بنموذج "المرونة المؤَمَّنة" (Flexicurité)، فإن الممارسة الميدانية أفرزت اختلالات بنيوية عميقة. وتتجسد الإشكالية المحورية لهذا الفصل في التساؤل التالي: \textbf{إلى أي حد استطاعت آليات الشغل المؤقت والمناولة التوفيق بين متطلبات التنافسية الاقتصادية للمقاولة وحماية المراكز القانونية والاجتماعية للأجراء، وما هي انعكاسات هذا النمط التعاقدي على دينامية المفاوضة الجماعية وهيكل سوق الشغل الوطني؟}

للإجابة عن هذه الإشكالية، يتناول هذا الفصل بالتحليل والتدقيق الآثار المترتبة عن هذه الأنماط على أربعة مستويات متكاملة:
\begin{itemize}[leftmargin=*, itemsep=4pt]
    \item \textbf{المستوى الأول}: تشخيص الانعكاسات المباشرة على الأجراء وحقوقهم المادية والمهنية؛
    \item \textbf{المستوى الثاني}: تقييم المكاسب والمخاطر الاستراتيجية والقانونية التي تواجه المقاولة؛
    \item \textbf{المستوى الثالث}: رصد تداعيات تشتت اليد العاملة على العمل النقابي ومؤسسات الحوار الاجتماعي؛
    \item \textbf{المستوى الرابع}: استقراء الآثار الهيكلية الشاملة على توازنات سوق الشغل وصناديق الحماية الاجتماعية.
\end{itemize}

%% ─────────────────────────────────────────────────────────
\section{الآثار على الأجراء}

يتحمل الأجير العبء الأكبر من تقلبات المرونة التعاقدية؛ إذ يجد نفسه في مركز قانوني واقتصادي متذبذب تتجاذبه التبعية المزدوجة بين مشغل قانوني (الوكالة أو مقاول المناولة) ومشغل فعلي يمارس الرقابة الميدانية اليومية (المقاولة المستعملة).

\subsection{تنامي الهشاشة الوظيفية وانعدام الأمان المهني}
تتجلى الهشاشة في قطاعي الشغل المؤقت والمناولة عبر مظاهر متعددة تؤثر على الاستقرار النفسي والمعيشي للأجير:
\begin{itemize}[leftmargin=*, itemsep=6pt]
    \item \textbf{انعدام الأفق التعاقدي المستمر}: تقتصر عقود الشغل المؤقت على مدة محددة حصرية بحكم المادة 500 من مدونة الشغل، مما يجعل الأجير يعيش في حالة قلق دائم وترقب مستمر لإنهاء الخدمة أو عدم تجديد عقد المهمة (Contrat de mission).
    \item \textbf{التقطع في الدخل والتبعية الاقتصادية}: تعاقب فترات العمل مع فترات بطالة غير معوضة (Inactivité intermittente) يحرم الأجراء من ضمان تدفق مالي مستقر، مما يؤدي إلى عجزهم عن الالتزام بالقروض البنكية أو الولوج للسكن الاقتصادي، ويكرس مفهوم "العامل الفقير" (Travailleur pauvre).
    \item \textbf{غياب المسار المهني والترقي الداخلي}: يتم استبعاد العمال المؤقتين وعمال المناولة من برامج التكوين المستمر والترقيات والسلالم الإدارية المتاحة للأجراء القارين داخل المؤسسة المستعملة، مما يجمد مهاراتهم في مستويات التنفيذ الدنيا.
\end{itemize}

\subsection{الإخلال بمبدأ المساواة في الأجور والامتيازات المادية}
أقر المشرع المغربي مبدأ المساواة في الأجر لمنع التمييز ضد الأجراء الموضوعين رهن الإشارة، وذلك بصريح العبارة في المادة 507 من مدونة الشغل:

\begin{tcolorbox}[colback=lightgray!30, colframe=maincolor, arc=3pt, title={\textbf{نص المادة 507 من القانون رقم 65.99 (مدونة الشغل)}}]
\textit{«يجب ألا يقل الأجر الذي يتقاضاه الأجير المؤقت عن الأجر الذي يتقاضاه أجير من نفس المؤهلات يشتغل بصفة دائمة لدى المقاولة المستعملة، وذلك بعد استيفاء فترة الاختبار.
يستفيد الأجير المؤقت من الامتيازات التي يستفيد منها أجراء المقاولة المستعملة في مجال النقل والمطاعم وغيرها من التسهيلات المماثلة...»}
\end{tcolorbox}

غير أن الواقع التطبيقي يكشف عن فجوة واسعة بين النص القانوني والممارسة العملية:
\begin{itemize}[leftmargin=*, itemsep=5pt]
    \item حرمان العمال المؤقتين من منح الأقدمية، ومنح المردودية السنوية (Primes de fin d'année / Treizième mois)، والتعويضات عن الأخطار الخاصة التي يستفيد منها نظراؤهم القارون.
    \item اللجوء في بعض قطاعات المناولة إلى تحديد الأجر الصافي بما يعادل الحد الأدنى القانوني حصراً (SMIG)، وتجريد الأجير من أي تعويض عن الساعات الإضافية الفعلية.
\end{itemize}

\subsection{تراجع التغطية الصحية ومخاطر حوادث الشغل والأمراض المهنية}
يُعدّ عمال المناولة والتشغيل المؤقت الشريحة الأكثر تعرضاً للإصابات الجسدية والأخطار البيئية:
\begin{itemize}[leftmargin=*, itemsep=6pt]
    \item \textbf{ضعف التكوين الوقائي الأولي}: إسناد مهام ذات طبيعة شاقة أو خطرة للأجراء دون إخضاعهم لتأهيل مسبق حول كيفية التعامل مع المواد الكيماوية أو الآلات الصناعية، بالمخالفة لمقتضيات المادة 281 من مدونة الشغل.
    \item \textbf{تشتت المسؤولية عند وقوع الحادثة}: تطبيقاً لمقتضيات القانون رقم 18.12 المتعلق بالتعويض عن حوادث الشغل، تقع المسؤولية المباشرة للتأمين على مشغل الأجير القانوني، غير أن المقاولات المستعملة غالباً ما تتنصل من توفير بيئة عمل آمنة معتبرة أن كلفة السلامة تقع على عاتق الشركة الموردة لليد العاملة.
    \item \textbf{عدم استيفاء شروط الاستفادة من الصندوق الوطني للضمان الاجتماعي (CNSS)}: يؤدي التقطع الزمني في عقود الشغل إلى عدم استكمال عدد الأيام القانونية المصرح بها (المشروطة في 3240 يوماً لاستحقاق معاش التقاعد، أو 54 يوماً خلال 6 أشهر للتغطية الصحية AMO)، مما يحرم عائلاتهم من التعويضات العائلية وخدمات العلاج المجاني.
\end{itemize}

%% ─────────────────────────────────────────────────────────
%% جدول مقارن لحقوق الأجراء
%% ─────────────────────────────────────────────────────────
\begin{table}[htbp]
\centering
\small
\renewcommand{\arraystretch}{1.5}
\begin{tabular}{|p{3.5cm}|p{5.5cm}|p{5.5cm}|}
\hline
\rowcolor{maincolor!20}
\textbf{معيار المقارنة} & \textbf{الأجير القار (عقد غير محدد المدة CDI)} & \textbf{الأجير المؤقت وعامل المناولة} \\
\hline
\textbf{الاستقرار الوظيفي} & أمان مهني محمي بمسطرة الفصل التأديبي والتعويض عن الفصل التعسفي. & هشاشة مستمرة، انتهاء العقد بانتهاء مدة المهمة دون تعويض عن الإعفاء. \\
\hline
\textbf{الأجر والتعويضات} & أجر أساسي + ترقيات + منح دورية وعلاوات إنتاج سنوية. & أجر مقترن بالحد الأدنى القانوني، مع حرمان متكرر من العلاوات التكميلية. \\
\hline
\textbf{الحماية الصحية والسلامة} & فحوصات طبية دورية منتظمة لدى طبيب الشغل، وتدريب مستمر على السلامة. & تدريب وقائي محدود، ومستوى تعرض مضاعف للأخطار المهنية داخل الورش. \\
\hline
\textbf{التقاعد والضمان الاجتماعي} & اشتراكات تراكمية مستمرة تضمن تقاعداً مريحاً وتغطية صحية شاملة. & فترات انقطاع متكررة تعيق تراكم الأيام القانونية لاستحقاق معاش التقاعد. \\
\hline
\end{tabular}
\vspace{0.6em}
\caption{مقارنة تحليلية للوضعية القانونية والمادية بين الأجراء الدائمين وأجراء الوساطة والمناولة}
\label{tab:comparaison-salaries}
\end{table}

%% ─────────────────────────────────────────────────────────
\section{الآثار على المقاولة}

تجد المقاولات الحديثة في آليات الشغل المؤقت والمناولة أداة بالغة الفعالية لإعادة هندسة مساطرها الإنتاجية، إلا أن التوسع المفرط فيها يُولّد تبعات استراتيجية وقانونية غير محمودة العواقب.

\subsection{المزايا التنافسية والمكتسبات الاقتصادية للمقاولة}
تسعى المقاولة المستعملة من وراء الاستعانة بشركات التشغيل المؤقت والمناولة إلى تحقيق مكاسب نوعية:
\begin{enumerate}[leftmargin=*, itemsep=6pt]
    \item \textbf{المرونة التسييرية والعددية (Flexibilité quantitative)}: القدرة اللحظية على مواءمة حجم قوة العمل مع دورات الإنتاج الصاعدة والنازلة وطلبيات التصدير الطارئة دون الخضوع للمساطر المعقدة لتقليص عدد الأجراء المنصوص عليها في المادة 66 وما يليها من مدونة الشغل.
    \item \textbf{التحكم في كتلة الأجور والتكاليف الثابتة}: تحويل جزء كبير من نفقات المستخدمين الثابتة (Frais de personnel) إلى خانة نفقات التسيير والخدمات الخارجية المتغيرة (Achats de services et charges externes)، مما يُحسّن من المؤشرات المالية المباشرة للمقاولة أمام الشركاء والمستثمرين.
    \item \textbf{التركيز على الكفاءات الجوهرية (Core Competencies)}: إسناد الأنشطة المساندة (كالخدمات اللوجستيكية، الصيانة التقنية، الحراسة، والنظافة) إلى مقاولات متخصصة ذات كفاءة تقنية أعلى وتكلفة تشغيلية أقل، مما يتيح للإدارة العامة تركيز مواردها على تطوير خطوط الإنتاج والابتكار التجاري.
\end{enumerate}

\subsection{المخاطر التنظيمية وفقدان الرأسمال المعرفي (Know-How)}
في المقابل، تدفع المقاولة ثمناً تنظيمياً غير مباشر ناتجاً عن تغييب الاستقرار البشري:
\begin{itemize}[leftmargin=*, itemsep=5pt]
    \item \textbf{نزيف الكفاءات وضياع المعرفة التراكمية}: يؤدي الدوران السريع للأجراء المؤقتين إلى عدم استقرار الخبرة العملية داخل فرق العمل، وصعوبة بناء ثقافة مؤسسية مشتركة قائمة على الثقة والولاء.
    \item \textbf{انخفاض الإنتاجية وجودة المنتجات}: تشير الدراسات الصناعية إلى أن غياب الحافز المهني لدى العامل المؤقت، وشعوره بالاغتراب المؤسسي داخل المعمل، ينعكسان مباشرة على ارتفاع معدلات الخطأ ونسب التلف في المواد الأولية.
    \item \textbf{تكاليف إعادة التأهيل المستمرة}: استنزاف وقت وجهد المشرفين ورؤساء الورشات في تدريب أجراء جدد مع بداية كل دورة تشغيل مؤقتة، مما يعطل وتيرة الإنتاج الاعتيادية.
\end{itemize}

\subsection{المسؤولية القانونية التضامنية والتبعات القضائية}
لم يُطلق المشرع يد المقاولة المستعملة في استغلال المرونة التعاقدية دون تحميلها مسؤوليات رادعة، بل قرر قواعد صارمة للمسؤولية التضامنية لحماية اليد العاملة:

\begin{tcolorbox}[colback=lightgray!30, colframe=maincolor, arc=3pt, title={\textbf{المسؤولية التضامنية للمقاولة الأصلية (المادتان 89 و90 من مدونة الشغل)}}]
\textit{«إذا كان المقاول من الباطن غير مقيد في السجل التجاري وغير مالك لأصل تجاري... فإن رئيس المقاولة الأصلية يتحمل المسؤولية في حالة عدم أداء المقاول من الباطن للأجور والمستحقات والتعويضات المستحقة لأجرائه.
كما تكون المقاولة الأصلية مسؤولة بالتضامن عن أداء اشتراكات الصندوق الوطني للضمان الاجتماعي في حالة عجز المقاول من الباطن...»}
\end{tcolorbox}

تتعرض المقاولة المستعملة والأصلية لمخاطر جسيمة عند عدم التدقيق في وضعية شركائها:
\begin{itemize}[leftmargin=*, itemsep=4pt]
    \item إمكانية الحكم عليها بالأداء المالي المباشر للأجور المتأخرة والتعويضات الناتجة عن حوادث الشغل.
    \item خطر إعادة تكييف العقود قضائياً (Requalification)؛ فإذا ثبت للقضاء أن المهمة المسندة للأجير المؤقت تدخل ضمن النشاط القار والعادي للمقاولة المستعملة، أو تجاوزت المدد القانونية، يُكيّف العقد على أنه عقد عمل دائم غير محدد المدة (CDI) تحت طائلة الحكم بتعويضات الفصل التعسفي.
\end{itemize}

%% ─────────────────────────────────────────────────────────
%% رسم بياني لمعادلة الآثار على المقاولة
%% ─────────────────────────────────────────────────────────
\vspace{0.8em}
\begin{center}
\begin{minipage}{0.92\textwidth}

\begin{tcolorbox}[colback=lightgray!30, colframe=maincolor, arc=4pt, width=\textwidth, boxrule=1pt, halign=center, valign=center]
\textbf{المدخلات: الاعتماد على التشغيل المؤقت والمناولة}\\[3pt]
\small خفض كلفة الأجور الثابتة، مرونة التكيف مع تقلبات الطلب، وتفويض المهام الثانوية.
\end{tcolorbox}

\centering{\color{maincolor}\Large$\Downarrow$}

\begin{tcolorbox}[colback=lightgray!30, colframe=maincolor, arc=4pt, width=\textwidth, boxrule=1pt, halign=center, valign=center]
\textbf{الآثار التنظيمية: تراجع الولاء والدوران الوظيفي المرتفع}\\[3pt]
\small فقدان المعرفة المهنية المتراكمة، ضعف الاندماج في فرق العمل، وتذبذب معايير الجودة.
\end{tcolorbox}

\centering{\color{maincolor}\Large$\Downarrow$}

\begin{tcolorbox}[colback=maincolor!15, colframe=maincolor, arc=4pt, width=\textwidth, boxrule=1.2pt, halign=center, valign=center]
\textbf{المخاطر القانونية: المسؤولية التضامنية وإعادة التكييف القضائي}\\[3pt]
\small تحمل أداء مستحقات الأجور والتغطية الاجتماعية ومواجهة دعاوى الفصل التعسفي.
\end{tcolorbox}

\vspace{0.8em}
\refstepcounter{figure}\label{fig:impact-entreprise}
\centering{\small\textbf{شكل \thefigure: دورة المخاطر والآثار التنظيمية للمرونة التعاقدية على المقاولة}}

\end{minipage}
\end{center}
\vspace{0.5em}

%% ─────────────────────────────────────────────────────────
\section{الآثار على الحوار الاجتماعي والعمل النقابي}

يُمثل توسع أنماط العمل الثلاثية أحد أكبر التحديات المعاصرة التي تواجه العمل النقابي والمفاوضة الجماعية بالمغرب، نظراً لما تفرضه من تفتيت لوحدة المصالح العمالية.

\subsection{تشتيت الكتلة العمالية وتفكيك وحدة الهيئة الناخبة}
تعتمد القوة التفاوضية للنقابات العمالية تاريخياً على التجانس المهني والتكتل الجماعي للأجراء داخل فضاء المنشأة الواحدة. إلا أن إدخال أجراء المناولة والعمال المؤقتين يخلق فصلاً مصطنعاً بين فئتين:
\begin{itemize}[leftmargin=*, itemsep=5pt]
    \item فئة "الأجراء الداخليين المحميين" (Insiders) الذين يتمتعون بعقود دائمة وأجور مستقرة وتمثيلية نقابية نافذة.
    \item فئة "الأجراء الخارجيين الهشين" (Outsiders) الذين يفتقرون للحق في التصويت في انتخابات مندوبي الأجراء بالمقاولة المستعملة، طبقاً للمادة 439 من مدونة الشغل التي تشترط أقدمية استقرار داخل المؤسسة لا تتوفر في الأجير المؤقت.
\end{itemize}
هذا التمايز يُفرغ التضامن العمالي من محتواه، ويحول دون توحيد المطالب الجماعية، بل قد يُولد صراعات وتنافساً غير صحي بين العمال أنفسهم للحفاظ على مناصبهم.

\subsection{عقبات ممارسة الحق النقابي والردع الضمني}
يكفل الدستور المغربي لسنة 2011 (الفصل 29) والاتفاقيات الدولية حق الانتماء والنشاط النقابي، غير أن الأجراء المؤقتين يصطدمون بموانع واقعية تحول دون ممارسة هذا الحق:
\begin{itemize}[leftmargin=*, itemsep=4pt]
    \item الهاجس المستمر من عدم تجديد عقد المهمة في حال الإقدام على الانخراط في النقابة أو المطالبة بالحقوق، مما يجعل مجرد التفكير في النشاط النقابي مرادفاً لفقدان مصدر الرزق.
    \item غياب مقرات تواصل نقابي خاصة بالأجراء المؤقتين داخل مقاولات التشغيل المؤقت نفسها، التي تكتفي بالتدبير الإداري للعقود دون توفير فضاءات للعمل الجمعي أو النقابي.
\end{itemize}

\subsection{تعقيد مسطرة المفاوضة الجماعية واتفاقيات الشغل}
تجد النقابات صعوبة قصوى في إدراج أجراء المناولة والتشغيل المؤقت في المفاوضة الجماعية (المواد 92 إلى 103 من مدونة الشغل) بسبب عائق ازدواجية رب العمل:
\begin{itemize}[leftmargin=*, itemsep=5pt]
    \item المقاولة المستعملة ترفض التفاوض بشأن تحسين أجورهم وظروف عملهم بدعوى انتفاء الرابطة التعاقدية المباشرة معهم، وتحيلهم على مشغلهم الأصلي.
    \item مقاولة التشغيل المؤقت ترفض بدورها الاستجابة للمطالب المالية بذريعة أن الفاتورة والتعريفة خاضعة حصراً للعقد التجاري المبرم مع المقاولة المستعملة، مما يضع الأجير في حلقة مفرغة تُجهض أي حوار اجتماعي مجدٍ.
\end{itemize}

%% ─────────────────────────────────────────────────────────
\section{الآثار على سوق الشغل}

يترك الانتشار الواسع لعقود المناولة والتشغيل المؤقت بصمات بالغة الوضوح على الهيكل العام لسوق الشغل والاقتصاد الوطني بالمغرب، موزعة بين أدوار إيجابية في التوظيف ومخاطر اجتماعية مقلقة.

\subsection{تسهيل الاندماج الأولي وتقليص كلفة البحث عن عمل}
تضطلع وكالات التشغيل المؤقت بدور الوسيط التفاعلي في التقريب بين العرض والطلب داخل سوق الشغل:
\begin{itemize}[leftmargin=*, itemsep=6pt]
    \item \textbf{جسر للولوج المهني للشباب}: توفر هذه الوكالات فرصة حيوية لحاملي الشواهد الجدد لاكتساب تجربة مهنية أولية تكسر حاجز اشتراط الخبرة السابقة لدى المقاولات الكبرى.
    \item \textbf{الحد من التكاليف المتبادلة للتوظيف}: تمكين المشغلين من تقييم كفاءة وملاءمة الأجير ميدانياً خلال فترة المهمة المؤقتة قبل تثبيته، مما يقلل من نسب الخطأ في التعيينات النهائية.
\end{itemize}

\subsection{الامتصاص التدريجي للقطاع غير المهيكل }
يساهم تقنين أنشطة المناولة والتشغيل المؤقت في دفع عجلة الحكامة في سوق الشغل عبر:
\begin{itemize}[leftmargin=*, itemsep=5pt]
    \item إلزام وكالات التشغيل بمسك سجلات قانونية ومحاسباتية رسمية وإيداع كفالات مالية والتقييد في السجل التجاري (المادتان 495 و496).
    \item تسجيل مئات الآلاف من العمال في أنظمة الضمان الاجتماعي والتأمين الإجباري الأساسي عن المرض (AMO)، والحد من التوظيف غير المصرح به (Travail au noir) السائد في قطاعات كالبناء والفلاحة والخدمات.
\end{itemize}

\subsection{المخاطر الهيكلية: استدامة الصناديق الاجتماعية وتوسع الهشاشة}
رغم مزايا الإدماج، فإن ترسيخ ثقافة الشغل المؤقت على حساب التوظيف الدائم يُحدث اختلالات هيكلية طويلة الأجل:
\begin{itemize}[leftmargin=*, itemsep=6pt]
    \item \textbf{تهديد التوازنات المالية لصناديق التقاعد}: يرتكز نظام المعاشات المدنية والمهنية على مبدأ التوزيع التضامني للأجيال، وهو ما يتطلب تدفقاً مستمراً للاشتراكات على مدار العام؛ ويؤدي تقطع الفترات المهنية إلى تراجع احتياطيات هذه الصناديق وتفاقم عجزها البنيوي.
    \item \textbf{ضعف القدرة الشرائية وركود الطلب الداخلي}: ترتبط وتيرة النمو الاقتصادي الداخلي باستقرار المداخيل العائلية المحفزة للاستهلاك والاستثمار الأسري طويل الأجل؛ وتؤدي الهشاشة المعيشية الناجمة عن العمل المؤقت إلى تراجع الإنفاق العائلي وتنامي القروض الاستهلاكية المتعثرة.
    \item \textbf{خطر تحول "الاستثناء المؤقت" إلى "قاعدة عامة"}: انزلاق سوق الشغل نحو نموذج يفتقر للعدالة الاجتماعية، حيث تتنصل المنشآت من الاستثمار في رأس المال البشري طويل الأجل، ما يُضعف إنتاجية العمل التراكمية على الصعيد الوطني.
\end{itemize}

%% ─────────────────────────────────────────────────────────
%% جدول تركيبي للآثار الإجمالية
%% ─────────────────────────────────────────────────────────
\begin{table}[htbp]
\centering
\small
\renewcommand{\arraystretch}{1.5}
\begin{tabular}{|p{3cm}|p{6cm}|p{6cm}|}
\hline
\rowcolor{maincolor!20}
\textbf{الطرف المتأثر} & \textbf{الآثار الإيجابية والمكاسب} & \textbf{الآثار السلبية والتحديات} \\
\hline
\textbf{الأجراء} & - اكتساب خبرة مهنية أولية.\newline - مدخل للتشغيل السريع وتفادي البطالة التامة. & - هشاشة وظيفية وغياب الأمان المهني.\newline - حرمان من بعض الامتيازات والمكافآت.\newline - ضعف التغطية الاجتماعية التراكمية. \\
\hline
\textbf{المقاولة} & - مرونة عددية فورية في تدبير الإنتاج.\newline - خفض النفقات الثابتة للأجور.\newline - التركيز على التخصص الأساسي للشركة. & - دوران وظيفي مرتفع وضياع المعرفة التقنية.\newline - تراجع الولاء للمؤسسة وجودة العمل.\newline - مخاطر المسؤولية التضامنية المالية والقانونية. \\
\hline
\textbf{الحوار الاجتماعي} & - إمكانية فتح قنوات حوار جديدة متخصصة مع وكالات التشغيل المؤقت وممثليها. & - تشتيت الكتلة العمالية والتمثيلية النقابية.\newline - صعوبة التفاوض الجماعي بسبب ازدواجية المشغل.\newline - عزوف العمال المؤقتين عن العمل النقابي. \\
\hline
\textbf{سوق الشغل} & - تسهيل التوفيق بين العرض والطلب.\newline - استقطاب جزء من القطاع غير المهيكل.\newline - تعزيز تنافسية الاستثمارات الخارجية. & - اتساع رقعة العمل الهش غير المستقر.\newline - إضعاف المداخيل الدورية لصناديق التقاعد.\newline - تذبذب مؤشرات الاستهلاك والقدرة الشرائية. \\
\hline
\end{tabular}
\vspace{0.6em}
\caption{حصيلة تركيبية لآثار التشغيل المؤقت والمناولة على المنظومة الاجتماعية والاقتصادية}
\label{tab:synthese-impacts}
\end{table}

%% ─────────────────────────────────────────────────────────
%% خاتمة الفصل
%% ─────────────────────────────────────────────────────────
\section*{خاتمة الفصل}
\addcontentsline{toc}{section}{خاتمة الفصل}

يتضح من خلال التحليل الميداني والمفاهيمي لآثار الشغل المؤقت والمناولة أن هذه الآليات التعاقدية لا يمكن اختزالها في مجرد أدوات تقنية لتدبير الموارد البشرية، بل هي خيار اقتصادي واجتماعي ذو أبعاد متعددة ونتائج متناقضة:

\begin{tcolorbox}[colback=lightgray!30, colframe=maincolor, arc=4pt, title={\textbf{خلاصة تركيبية للفصل السابع}}]
\begin{itemize}[leftmargin=*, itemsep=4pt]
    \item \textbf{ازدواجية الأثر}: تمنح هذه الأنماط المقاولة متنفساً حيوياً لتحقيق المرونة التنافسية وامتصاص صدمات الأسواق، وتفتح للأجراء الجدد آفاق الولوج لسوق العمل، إلا أنها تفرض كلفة اجتماعية باهظة تُلقي بظلالها على الاستقرار المادي والنفسي للأجراء.
    \item \textbf{تآكل التضامن المؤسسي}: أفرزت سيولة اليد العاملة تراجعاً ملحوظاً في فاعلية العمل النقابي والمفاوضة الجماعية، وخلقت انقساماً بنيوياً داخل المنشأة الواحدة بين أجراء مستقرين يتمتعون بالحقوق الكاملة، وأجراء مؤقتين يعيشون على هامش الحماية المهنية.
    \item \textbf{الحاجة إلى إعادة الضبط التشريعي}: إن مواجهة مخاطر هشاشة العمل واستنزاف صناديق الحماية الاجتماعية لا تقتضي منع هذه الأنماط المرنة، بل تستوجب تحيين الإطار القانوني وتعزيز منظومة الرقابة التفتيشية لضمان احترام مبدأ المساواة في الأجر وتفعيل المسؤولية التضامنية للمقاولات المستعملة.
\end{itemize}
\end{tcolorbox}

تؤسس هذه الخلاصات الاجتماعية والاقتصادية بشكل مباشر لأهمية الانتقال إلى دراسة الواقع العملي التطبيقي؛ حيث ستتيح \textbf{دراسة الحالة الواقعية} في الفصل الموالي تفكيك هذه الآثار على أرض الواقع من خلال وضعية مهنية حية، ومعاينة أدوار مختلف المتدخلين (مدير المصنع، مسؤولو المناولة، ممثلو الأجراء، ومفتش الشغل) في تدبير هذه التوازنات المعقدة.